% Urban Imperviousness Mapping with AlphaEarth Embeddings and Sentinel-2:
% Local Performance, Cross-City Transfer and Training-Label Effects.
%
% Separate journal-style paper, distinct from report/IMD_Mapping_report.tex.
% Outline from Prof. Brovelli, saved at paper/OUTLINE.md. Written one section
% at a time; this file currently holds Section 1 (Introduction) only.
%
% Compile with XeLaTeX: `latexmk` (see latexmkrc). This directory is
% self-contained: paper.tex, refs.bib and figs/ (once added) are everything
% the build needs.
\documentclass[11pt,a4paper]{article}

\usepackage[margin=2cm]{geometry}
\usepackage{fontspec}
\setmainfont{Times New Roman}[Ligatures=TeX]

\usepackage{booktabs}
\usepackage{amsmath}
\usepackage{siunitx}
\sisetup{group-separator={,}, group-minimum-digits=4}
\usepackage{graphicx}
\usepackage[font=small,labelfont=bf,justification=justified,
            singlelinecheck=false,skip=6pt]{caption}
\usepackage{array}
\usepackage{multirow}
\usepackage{float}
\usepackage{xcolor}
% Layout for paper.tex, \input from its preamble right after xcolor.
% Copied from report/layout.tex -- same typesetting conventions, same project.

% Clickable links (URLs, citations, cross-references).
\definecolor{linkcol}{HTML}{467886}

% Float placement.
%
% Kept in its own file because these are typesetting parameters, not claims:
% an audit script that reads paper.tex for numbers would otherwise report
% every fraction below as an unsourced value. Nothing here is a result.
\renewcommand{\topfraction}{0.9}
\renewcommand{\bottomfraction}{0.8}
\renewcommand{\textfraction}{0.07}
\renewcommand{\floatpagefraction}{0.8}
\setcounter{topnumber}{3}
\setcounter{bottomnumber}{2}
\setcounter{totalnumber}{5}
\setlength{\floatsep}{8pt plus 2pt minus 2pt}
\setlength{\textfloatsep}{10pt plus 2pt minus 4pt}
\setlength{\intextsep}{8pt plus 2pt minus 2pt}

% Two floats side by side in one figure environment: half the text width
% less a gutter.
\newlength{\halfwidth}
\setlength{\halfwidth}{\dimexpr0.48\textwidth\relax}

\usepackage[numbers,square,sort&compress]{natbib}
\bibliographystyle{unsrtnat}
\renewcommand{\bibfont}{\small}
\usepackage[colorlinks,breaklinks,
            linkcolor=linkcol,citecolor=linkcol,urlcolor=linkcol]{hyperref}
\usepackage{xurl}
\usepackage[capitalise,nameinlink]{cleveref}

\graphicspath{{figs/}}

\newcommand{\Rsq}{$R^2$}

\setlength{\parindent}{0pt}
\setlength{\parskip}{6pt}

% Author list copied from the internal report as a default -- not yet
% confirmed by Prof. Brovelli for this paper specifically. Update once the
% author list and order for submission are settled.
\title{Urban Imperviousness Mapping with AlphaEarth Embeddings and
Sentinel-2: Local Performance, Cross-City Transfer and Training-Label Effects}
\author{Maria Antonia Brovelli, Matej Žgela, Keerthana Kirubakaran and Xiao Tan\\[2pt]
\small Politecnico di Milano}
\date{}

\begin{document}
\maketitle

\section{Introduction}
\label{sec:intro}

\subsection{Imperviousness density and applications}
\label{sec:intro-imd}

Imperviousness, or impervious surface density (IMD), is the proportion of
land surface covered by impervious materials such as buildings, roads and
pavements. Unlike natural ground, these surfaces do not let rain infiltrate,
and turn it into runoff instead. IMD mapping estimates the impervious
fraction for each pixel of a study area and produces a continuous spatial
map, with every pixel holding a value from 0 to 100\%, rather than a binary
built/not-built label.

Because it is continuous and spatially explicit, IMD underpins a range of
applications beyond the surface-cover inventory it directly measures.
Stormwater and flood-risk models use it to estimate runoff generation at the
catchment scale. Urban Heat Island (UHI) studies use it as a covariate for
land surface temperature, since impervious materials retain and re-radiate
heat differently from vegetated or bare ground. Local Climate Zone (LCZ)
classification, which groups urban areas into classes with comparable
thermal behaviour, uses IMD as one of its defining surface-cover fractions.
Sustainability and land-monitoring programmes track IMD change over time as
an indicator of urban growth and soil sealing.

This range of downstream uses is also what motivates the present study. Milan
has a published IMD product; Hanoi and Ho Chi Minh City (HCMC) do not. A map
is needed for both cities to support LCZ mapping and UHI analysis within the
Italy--Vietnam LCZ-UHI-GEO project and the \emph{Space it up!} project, whose
partners are Politecnico di Milano and the Vietnam National Space Center
(VNSC). Producing it is the applied motivation behind the comparison this
paper reports.

\subsection{Existing EO-based IMD mapping: state of the art}
\label{sec:intro-sota}

Mapping impervious surfaces from Earth observation (EO) imagery has been an
active area of remote sensing research for over two decades; Weng
\citep{weng2012impervious} reviews the methods, requirements and trends up to
the early 2010s. Two broad families of approach dominate the literature.
Sub-pixel methods, most commonly spectral mixture analysis, decompose each
pixel's reflectance into fractions of a small set of pure materials
(impervious surface, vegetation, soil, sometimes water) and report the
impervious fraction directly. Regression and classification methods instead
train a statistical or machine-learning model -- historically linear
regression and support vector machines, more recently ensemble
tree methods -- against a reference IMD or built-up surface layer, then apply
it to new imagery. The present study follows the second family: a
Random Forest regressor trained against a reference product (\cref{sec:data}),
the approach also used by the two reference products this study relies on
(\cref{sec:data}).

Global and regional IMD or built-up-surface products exist at a range of
resolutions. The Copernicus Land Monitoring Service's Imperviousness Density
and the Global Human Settlement Layer's GHS-BUILT-S, both used as training
targets in this study (\cref{sec:data}), are two examples. Others combine
multiple sensors on cloud platforms to reach global or near-global coverage:
Zhang et al.\ \citep{zhang2020global30m} produce a 30\,m global impervious
surface map for 2015 by combining Landsat~8, Sentinel-1 and VIIRS night-light
imagery on Google Earth Engine, and Sun et al.\ \citep{sun2023mekong} map
10\,m impervious surface for the Greater Mekong Subregion -- which includes
northern and southern Vietnam -- from Sentinel-2 and Sentinel-1. Neither,
however, produces a 2018-dated, 10\,m product for Hanoi or HCMC specifically,
which is the gap this study's reference-target choice (\cref{sec:data}) works
around rather than fills directly.

A recurring finding across this literature is that mapping accuracy depends
as much on the reference product a model is trained against as on the
predictor imagery itself, since different reference products encode
different physical definitions of what counts as impervious (sealed ground
versus roofed structure, for instance). \Cref{sec:results} returns to this
point directly for the three cities studied here.

\subsection{Geospatial foundation models and embeddings}
\label{sec:intro-fm}

Foundation models -- large models pretrained on broad data and adapted to
many downstream tasks -- have more recently been extended from natural
imagery to Earth observation. Xiao et al.\ \citep{xiao2024fmsurvey} survey
this emerging field, covering visual, vision-language and other remote
sensing foundation models and the growing collection of pretraining datasets
built to support them.

One strand of this work produces \emph{embeddings}: fixed-length, per-pixel
feature vectors that summarise large amounts of multi-sensor, multi-temporal
satellite imagery into a single, ready-made input layer. Used this way, a
downstream task needs no local image compositing -- no choice of which
dates to use or how to combine them -- because that work has already been
done once, globally, by the model that produced the embeddings. Google's
AlphaEarth Foundations \citep{brown2025alphaearth} is one such model,
providing 64-dimensional annual embeddings worldwide through Google Earth
Engine \citep{gorelick2017gee}.

Most published evaluations of EO embeddings address discrete tasks -- land
cover classification, crop type mapping -- rather than continuous surface
estimation. Whether embeddings of this kind are competitive with purpose-built
local composites for a continuous regression target like IMD, and whether
a model trained on them transfers across cities with different climates and
urban forms, is accordingly still largely untested. This is the gap the
present study addresses.

\subsection{Research questions and contributions}
\label{sec:intro-rq}

This paper asks three questions, corresponding to the three themes in its
title.

\begin{enumerate}
  \item \textbf{Local performance.} For 10\,m imperviousness mapping within a
        single city, how does a Random Forest trained on AlphaEarth
        embeddings compare with one trained on a locally composed
        Sentinel-2 image, under the same sampling, validation and tuning
        procedure?
  \item \textbf{Cross-city transfer.} Does a model trained in one city apply
        directly to another, climatically and structurally different city
        (zero-shot transfer), or is local retraining needed -- and does the
        answer differ between AlphaEarth and Sentinel-2 predictors?
  \item \textbf{Training-label effects.} How much of a mapping method's
        accuracy is attributable to the predictor (AlphaEarth versus
        Sentinel-2) as opposed to the reference product it is trained
        against (CLMS versus GHS-BUILT-S)?
\end{enumerate}

Answering them, the paper makes four contributions. First, it compares
Sentinel-2 and AlphaEarth embeddings for local 10\,m imperviousness mapping
in Milan, under an identical spatial-block sampling, cross-validation and
hyperparameter-tuning procedure for both predictors. Second, it evaluates
both zero-shot transfer and local retraining of that mapping across three
cities with different climates and urban morphologies -- Milan, Hanoi and
HCMC -- none of which has previously been compared this way for this task.
Third, it isolates the effect of the training target by producing Milan maps
against both CLMS and GHS-BUILT-S under otherwise identical conditions.
Fourth, it validates every resulting map against an independent reference
dataset, EarthLabel, that was not used to train any of the models, so that
the comparisons above are not an artefact of a shared reference product.

\end{document}
